\documentclass[conference]{IEEEtran}
\IEEEoverridecommandlockouts
% The preceding line is only needed to identify funding in the first footnote. If that is unneeded, please comment it out.
\usepackage{cite}
\usepackage{amsmath,amssymb,amsfonts}
\usepackage{algorithmic}
\usepackage{graphicx}
\usepackage{textcomp}
\usepackage{xcolor}
%\usepackage{booktabs}
\usepackage{tabularray}
\usepackage{flushend}
\UseTblrLibrary{booktabs}
\def\BibTeX{{\rm B\kern-.05em{\sc i\kern-.025em b}\kern-.08em
    T\kern-.1667em\lower.7ex\hbox{E}\kern-.125emX}}
\begin{document}

\title{Experimenting with space constrained shortwave antennas}

\author{\IEEEauthorblockN{Gergely Vakulya}
\IEEEauthorblockA{\small \textit{Alba Regia Technical Faculty} \\
\textit{Óbuda University}\\
\textit{Székesfehérvár, Hungary}\\
\textit{vakulya.gergely@amk.uni-obuda.hu}
}}

\maketitle

\begin{abstract}
  Designing good antennas for a specific application is always
  a complicated problem. The different design goals (high gain,
  directional or uniform radiation pattern, easily matchable
  and stable impedance) are not easy to achieve due to
  physical constrains, most frequently the limited space.
  This gets even more problematic with lower frequencies,
  especially on the shortwave bands.

  This paper presents 4 portable shortwave antennas, which
  were designed for demonstration purposes. The important
  and often misunderstood properties of the antennas
  are summarized and presented.
\end{abstract}

\begin{IEEEkeywords}
antenna, shortwave
\end{IEEEkeywords}

\section{Introduction}


\section{Related work}

\begin{figure}[htbp]
	\centering
	\includegraphics[width=0.35\textwidth]{fig/dummy.png}
	\caption{Dummy image}
	\label{fig-dummy}
\end{figure}

\section{Summary}


\bibliographystyle{IEEEtran}
\bibliography{IEEEabrv,references}

\end{document}
